% Introduction.
Comparisons between artificial networks and neural data have become a standard way to ask which
models capture the computations of visual cortex. Representational similarity analysis (RSA) is one
of the most widely used tools for this purpose: it summarises each system by the dissimilarities
between its responses to a set of stimuli and compares these summaries directly, without fitting a
mapping \citep{kriegeskorte2008}. A recurring and surprising result is that networks with random,
untrained weights can rival trained networks as models of visual cortex, in some reports exceeding
them \citep{truzzi2020, kazemian2025}. If true, such results bear on how much of early visual
representation follows from architecture alone and how much from learning.

We reported a result of this kind ourselves. In THINGS fMRI \citep{hebart2023}, an untrained network
with uncalibrated batch-normalisation statistics exceeded a backpropagation-trained network at V1
when images were evaluated at 224\,px, but not at 32\,px, the resolution at which the networks had
been trained \citep{evalres}; a related claim appears in \citet{leutenegger2026}. Both analyses used
a single convention: globally pooled features, correlation distance between units, and voxel
responses z-scored over all trials of the experiment. This paper traces that claimed property of
untrained networks to a normalisation convention.

Correlation distance across units is not invariant to constant offsets of individual units over the
stimulus set: when units have different mean responses, the dissimilarity between two stimuli depends
in part on how strongly each pattern is aligned with the shared profile of unit means. The same holds
for voxels. Methodological work has shown that removing or keeping the mean of the measured patterns
changes what RSA measures: demeaning can produce ``representational confusion''
\citep{ramirez2017}, the regional-mean response profile carries information that standard RSA
discards \citep{taylor2025}, and averaging over neurons within voxels amplifies activity shared by
all sampled neurons, a distortion that mean removal largely corrects \citep{bossio2025}. Centring is
also explicit in related similarity measures \citep{kornblith2019}, and the choice of dissimilarity
measure affects the reliability of RSA \citep{walther2016, diedrichsen2017}. Whether such choices
reverse the ranking of candidate models has rarely been tested.

We make six contributions. (i)~A pre-specified ridge readout on the full feature map, in THINGS and,
with the same plan, in the Natural Scenes Dataset \citep[NSD;][]{allen2022}, places the trained network
ahead at every pooling grid. (ii)~A multiverse of 36 combinations of unit normalisation, distance
measure and voxel normalisation in THINGS \citep{steegen2016} shows that the sign of the V1 contrast
depends on the analysis. (iii)~A replication in NSD (8 subjects) finds no analysis in which the
untrained network is ahead under the NSD standard normalisation. (iv)~We identify the convention that
creates the advantage: a voxel response profile common to all analysed trials, which is present in unnormalised betas and, in
THINGS, also after normalising over all trials of the experiment, because all 720 images were shown in
one session; a pre-specified prediction that the advantage appears in NSD when this profile is left in
(raw betas) was confirmed. (v)~We describe the
mechanism as far as it is identified: with the profile left in, both representational dissimilarity
matrices (RDMs) are dominated by per-image response amplitude, and the untrained network's RDM more so.
(vi)~We give recommendations for model--brain comparisons. Every pre-specified prediction is reported,
including the \FailedPredictions{} that were not met, and we revise the corresponding claims of our
earlier work.
